\section{Программная реализация прототипа и модуля сравнения с эталоном}
\label{sec:razdel4}

\subsection{Разработка тестового Go-приложения как управляемой нагрузки}
\label{sec:4.1}

Пакет \texttt{app} (листинг полностью приведён в приложении~\ref{app:code})
реализует единственный тип данных \texttt{Record} (запись с полями
\texttt{ID}, \texttt{Name}, \texttt{Category}, \texttt{Amount}, \texttt{Active})
и шесть экспортируемых функций обработки срезов таких записей:
\texttt{ParseRecords} (разбор JSON), \texttt{FilterActive} (фильтрация по
булеву полю), \texttt{FindByID} (линейный поиск), \texttt{Aggregate}
(группировка сумм по категории через map), \texttt{FormatNames}
(форматирование в строки) и \texttt{Deduplicate} (удаление повторов по
\texttt{ID} с сохранением первого вхождения). Для каждой функции написан
бенчмарк (\texttt{app/\allowbreak bench\_test.go}) на синтетическом наборе из 1000
записей (\texttt{genRecords}), генерируемом детерминированно (без
использования \texttt{math/\allowbreak rand}), что обеспечивает воспроизводимость
измерений между прогонами. Для \texttt{FindByID} бенчмарк намеренно ищет
заведомо отсутствующий \texttt{ID}, чтобы гарантировать полный проход по
срезу на каждой итерации — иначе результат зависел бы от случайной позиции
искомого элемента.

Отдельно, в \texttt{app/\allowbreak processor\_test.go}, реализованы unit-тесты,
проверяющие корректность бизнес-логики независимо от производительности
(таблица~\ref{tab:4.1}) — эти тесты гарантируют, что скрипты сравнения
(раздел~\ref{sec:4.3}) в сценариях \texttt{regression/\allowbreak *}
(раздел~\ref{sec:razdel5}) сравнивают производительность функционально
эквивалентного кода, а не случайно сломанной реализации.

\begin{table}[htbp]
\caption{Unit-тесты пакета \texttt{app}}
\label{tab:4.1}
\centering
\begin{tabular}{ll}
\toprule
Тест & Проверяемое свойство \\
\midrule
\texttt{TestParseRecords} & корректный разбор, пустой массив, некорректный JSON \\
\texttt{TestFilterActive} / \texttt{\_Empty} & фильтрация по \texttt{Active}, включая пустой вход \\
\texttt{TestFindByID} & найденный и отсутствующий \texttt{ID} \\
\texttt{TestAggregate} & суммирование по категориям \\
\texttt{TestFormatNames} & формат строкового представления \\
\texttt{TestDeduplicate} & отсутствие повторов \texttt{ID} на выходе \\
\bottomrule
\end{tabular}
\end{table}

\subsection{Реализация критерия сравнения с эталоном}
\label{sec:4.3}

Критерий~\eqref{eq:criterion} реализован в \texttt{scripts/\allowbreak lib.sh} функцией
\texttt{check\_regression\_csv}, принимающей CSV-отчёт \texttt{benchstat} и
порог $T$ (листинг~\ref{lst:lib}). Вынесение этой функции в отдельный
переиспользуемый файл — прямое следствие вывода раздела~\ref{sec:3.3}: до
рефакторинга (см. историю коммитов прототипа) идентичная логика
дублировалась в \texttt{compare.sh} и в каждом из экспериментальных
скриптов раздела~\ref{sec:razdel5}, что создавало риск рассинхронизации
критерия между решающим и экспериментальными путями. Функция построчно
разбирает CSV, отслеживая текущую метрику (\texttt{sec/\allowbreak op}, \texttt{B/\allowbreak op},
\texttt{allocs/\allowbreak op} — секции CSV-отчёта \texttt{benchstat} размечены пустой
первой колонкой), отбирает строки с положительным изменением
(\texttt{vs\_base} начинается на \texttt{+}, поскольку все три метрики
lower-is-better) и через \texttt{awk} сравнивает величину изменения с
порогом.

\begin{lstlisting}[language=bash,caption={Критерий регрессии --- scripts/lib.sh (сокращено)},label={lst:lib}]
check_regression_csv() {
  local csv="$1" threshold="$2"
  local metric=""
  REGRESSED=0
  ROWS=""
  while IFS=, read -r name base_val base_ci cur_val cur_ci vs_base p_col; do
    if [[ "$name" == "" && "$base_val" == "sec/op" ]]; then metric="ns/op"; continue; fi
    # ... аналогично для B/op, allocs/op
    [[ "$name" == "" || "$name" == "geomean" ]] && continue
    [[ "${vs_base:0:1}" != "+" ]] && continue   # интересует только рост

    local magnitude="${vs_base#+}"; magnitude="${magnitude//%/}"
    if awk -v m="$magnitude" -v t="$threshold" 'BEGIN{exit !(m > t)}'; then
      REGRESSED=1
      ROWS="${ROWS}| ${metric} | ${name} | ${vs_base} | ... |"$'\n'
    fi
  done <"$csv"
}
\end{lstlisting}

Условие $p_{b,m} < \alpha$ из критерия~\eqref{eq:criterion} в явном виде в
листинге~\ref{lst:lib} не присутствует: оно обеспечивается тем, что
\texttt{benchstat} вызывается с флагом \texttt{-alpha}, и сам не печатает
относительное изменение (столбец \texttt{vs\_base}) для строк, где
отличие статистически не значимо на этом уровне — вместо процента в CSV
оказывается символ незначимости, который не начинается с \texttt{+} и
поэтому естественно отфильтровывается условием \texttt{[[
"\$\{vs\_base:0:1\}" != "+" ]]}. Такое распределение ответственности
(``arithметика значимости'' — в \texttt{benchstat}, ``арифметика порога'' —
в \texttt{lib.sh}) соответствует границе компонентов, установленной
архитектурой (рис.~\ref{fig:3.1}).

Скрипт \texttt{scripts/\allowbreak compare.sh} — точка входа решающего шага: запускает
\texttt{go test -bench=. -benchmem -count=\$COUNT}, передаёт результат и
\texttt{baseline.txt} в \texttt{benchstat -format=csv}, вызывает
\texttt{check\_regression\_csv} и завершает процесс с кодом \texttt{1} при
обнаруженной регрессии — код завершения интерпретируется GitHub Actions как
падение шага и, соответственно, блокирует required-check pull request'а.

Диагностический шаг реализован в \texttt{scripts/\allowbreak profile-diff.sh}: те же
бенчмарки перезапускаются с флагами \texttt{-cpuprofile}/\texttt{-memprofile},
после чего \texttt{go tool pprof -top -diff\_base=<эталон> -nodecount=10}
формирует top-10 функций по каждому из двух профилей (CPU по cumulative
time, heap по alloc\_space) и печатает результат в markdown-блоке кода —
пригодном для непосредственной вставки в \texttt{\$GITHUB\_STEP\_SUMMARY}
без дополнительного форматирования.

\subsection{Настройка CI/CD-пайплайна с профилированием}
\label{sec:4.2}

CI-пайплайн (\texttt{.github/\allowbreak workflows/\allowbreak perf-check.yml}, листинг~\ref{lst:ci})
запускается на каждое открытие/обновление pull request в \texttt{main} и
последовательно выполняет оба шага модели процесса (рис.~\ref{fig:2.1}):
решающий (\texttt{compare.sh}) и диагностический (\texttt{profile-diff.sh},
запущенный с \texttt{if: always()}, то есть выполняющийся даже если
предыдущий шаг завершился неудачей — иначе профиль для диагностики
регрессии был бы недоступен именно в том прогоне, где она обнаружена).
Сырые pprof-профили текущего прогона сохраняются шагом
\texttt{actions/\allowbreak upload-artifact@v4}~\cite{gh-upload-artifact} и доступны для скачивания из интерфейса
GitHub Actions в течение 14 дней — разработчик может открыть их локально
интерактивным профилировщиком (\texttt{go tool pprof -http=:8080
downloaded-cpu.pprof}) без повторного запуска бенчмарков на своей машине.

\begin{lstlisting}[language=bash,caption={CI-пайплайн --- .github/workflows/perf-check.yml (ключевые шаги)},label={lst:ci}]
on:
  pull_request:
    branches: [main]
jobs:
  benchmark:
    runs-on: ubuntu-latest
    steps:
      - uses: actions/checkout@v4
      - uses: actions/setup-go@v5
        with: { go-version: "1.22" }
      - name: Install benchstat
        run: go install golang.org/x/perf/cmd/benchstat@latest
      - name: Compare benchmarks against baseline
        run: ./scripts/compare.sh | tee "$GITHUB_STEP_SUMMARY"
      - name: Profile and diff against baseline
        if: always()
        env: { CUR_DIR: /tmp/pprof-out }
        run: ./scripts/profile-diff.sh | tee -a "$GITHUB_STEP_SUMMARY"
      - name: Upload pprof profiles
        if: always()
        uses: actions/upload-artifact@v4
        with: { name: pprof-profiles, path: /tmp/pprof-out/*.pprof }
\end{lstlisting}

Основные категории пользователей прототипа: \emph{разработчик} — автор
pull request'а, локально запускающий \texttt{compare.sh}/\texttt{profile-diff.sh}
перед отправкой изменений и читающий отчёт CI после отправки;
\emph{ревьюер} — использует статус required-check и markdown-отчёт как
основание для решения об одобрении слияния, не запуская инструменты
самостоятельно; \emph{сопровождающий пайплайна} — обновляет эталонные
данные (\texttt{scripts/\allowbreak capture-baseline-profiles.sh}, вручную по
необходимости) при осознанном изменении производительности эталонной
ветки. Основной сценарий использования — линейная последовательность:
разработчик открывает PR $\to$ CI выполняет оба шага $\to$ отчёт публикуется
в теле PR (рис.~\ref{fig:4.1}, пример реального вывода) $\to$ ревьюер
принимает решение на основе статуса проверки.

\begin{figure}[htbp]
\centering
\begin{lstlisting}[language={},basicstyle=\ttfamily\fontsize{8.5}{10}\selectfont,frame=single,backgroundcolor=\color{codebg}]
## Отчёт сравнения производительности

Порог: 10%, alpha: 0.05, повторов: 10, baseline: benchmarks/baseline.txt

| Метрика | Бенчмарк | Δ | p |
|---|---|---|---|
| ns/op | Deduplicate-8 | +62.33% | 0.000 |

FAIL: обнаружена регрессия производительности
\end{lstlisting}
\caption{Реальный вывод \texttt{compare.sh} на ветке regression/quadratic-dedup (величина отклонения варьируется между запусками из-за шума окружения раннера --- в независимых прогонах наблюдались значения от $+62\%$ до $+69\%$, см. раздел~\ref{sec:5.4})}
\label{fig:4.1}
\end{figure}

\subsection{Тестирование прототипа}
\label{sec:4.4}

Тестирование прототипа выполнялось на двух независимых уровнях. Во-первых,
корректность \emph{бизнес-логики} тестового приложения — стандартными
Go unit-тестами (таблица~\ref{tab:4.1}), выполняемыми командой \texttt{go
test ./\allowbreak ...} и не связанными с производительностью. Во-вторых, корректность
\emph{самого механизма обнаружения регрессий} — тест-планом, где каждая из
четырёх веток \texttt{regression/\allowbreak *} (описаны подробно в
разделе~\ref{sec:5.1}) выступает тестовым случаем с заранее известным
ожидаемым исходом ("должен провалиться"), а ветка \texttt{main} (без
внесённых изменений) — контрольным случаем с ожидаемым исходом "не должен
провалиться" (проверка отсутствия ложных срабатываний). Такой тест-план
формально эквивалентен приёмочному тестированию: вместо того чтобы
проверять выход функции на фиксированном входе, он проверяет выход всего
пайплайна (\texttt{compare.sh}) на фиксированном изменении кода.
Количественные результаты этого тестирования, включая перебор сетки
параметров $T \times \alpha$, вынесены в раздел~\ref{sec:razdel5}, поскольку
представляют самостоятельный экспериментальный результат, а не просто
факт прохождения/непрохождения теста.

Сравнение с аналогами (таблица~\ref{tab:1.1}) в части "достигнутых
системных показателей" сведено к двум характеристикам, для которых
у сопоставимых Go-ориентированных инструментов (benchmark-action,
CodSpeed) отсутствуют открыто публикуемые числа: (а) корректность
обнаружения на управляемых сценариях (раздел~\ref{sec:5.2}) и
(б) накладные расходы pprof-профилирования, добавленные к длительности
CI-джобы (раздел~\ref{sec:5.3}) — сопоставление именно этих двух
характеристик и составляет предмет экспериментальной апробации.

\subsection{Выводы}

\begin{enumerate}
\item Реализовано тестовое Go-приложение (пакет \texttt{app}, 6 функций,
6 бенчмарков, 6 unit-тестов) как управляемая, воспроизводимая нагрузка для
демонстрации работы системы профилирования; воспроизводимость обеспечена
детерминированной генерацией тестовых данных без использования
источников случайности.
\item Реализован критерий обнаружения регрессии~\eqref{eq:criterion} в виде
переиспользуемой функции \texttt{check\_regression\_csv}
(\texttt{scripts/\allowbreak lib.sh}), вызываемой как решающим скриптом
(\texttt{compare.sh}), так и всеми тремя экспериментальными скриптами
раздела~\ref{sec:razdel5} — это устраняет риск рассинхронизации логики
критерия между продуктивным и экспериментальным путями выполнения.
\item Реализован диагностический механизм \texttt{profile-diff.sh},
дополняющий решающий шаг: снимает CPU/heap pprof-профили текущего прогона
и сравнивает их с эталонными на уровне отдельных функций через
\texttt{go tool pprof -diff\_base}, публикуя top-10 функций по каждому
профилю в отчёте.
\item Оба шага интегрированы в единый CI-пайплайн на GitHub Actions,
выполняющийся на каждый pull request; диагностический шаг настроен на
выполнение независимо от исхода решающего (\texttt{if: always()}), а сырые
профили сохраняются как артефакт джобы на 14 дней для последующего
интерактивного анализа.
\item Разработан двухуровневый тест-план прототипа: unit-тесты бизнес-логики
приложения (не зависят от производительности) и приёмочный тест-план
механизма обнаружения регрессий на четырёх управляемых сценариях плюс
контрольном сценарии без изменений — количественные результаты вынесены
в раздел~\ref{sec:razdel5} как самостоятельный экспериментальный результат.
\end{enumerate}
